\chapter{Puissances et racines carrées}
\section*{Exercice 4.1 -- Puissances}
Les exposants donnent $4-7+1=-2$, donc $1/4$ ; puis $(-2)\times3+5=-1$, donc $1/3$. Les bases sont non nulles, ce qui autorise les exposants négatifs.

\section*{Exercice 4.2 -- Notation scientifique}
Le coefficient vaut $6\times4/3=8$ et l'exposant $-3+7-2=2$. Le résultat est $8\times10^2=800$.

\section*{Exercice 4.3 -- Racines}
La première expression vaut $5\sqrt2+4\sqrt2-3\sqrt2=6\sqrt2$. La seconde vaut $|-7|=7$, car une racine carrée est non négative.

\section*{Exercice 4.4 -- Comparer}
Les deux nombres du premier couple sont positifs ; leurs carrés valent $28$ et $25$, donc $2\sqrt7>5$. Pour le second, le carré de la somme vaut $9+2\sqrt{14}>9$, donc $\sqrt2+\sqrt7>3=\sqrt9$.

\section*{Exercice 4.5 -- Seuil}
$1{,}5^5=7{,}59375<10$ et $1{,}5^6=11{,}390625>10$. La suite des puissances est strictement croissante puisque le facteur est supérieur à $1$. Le premier rang est donc $6$.

\section*{Exercice 4.6 -- Une preuve}
Le quotient de racines est non négatif et son carré vaut $a/b$ ; l'unicité de la racine donne l'égalité. Pour $a=b=-1$, le quotient $a/b=1$ possède bien une racine, mais $\sqrt{-1}$ n'est pas un réel : le membre de droite n'existe pas dans $\R$.
